\documentclass{article}
\usepackage[pdftex]{graphicx}
\usepackage{xcolor}
\usepackage[letterpaper, margin = 2.5cm]{geometry}
\usepackage[T2A]{fontenc}
\usepackage[english, russian]{babel}
\usepackage{amsmath, amsfonts, amssymb}
\usepackage{hyperref, url}

\usepackage{minted}

\bibliography{refs.bib}

\usepackage{parskip}

\begin{document}

\title{Домашнее задание 1}
\date{}
\author{SMR-2026 Starkit Education, Илья Осокин @elijahmipt} %поменяйте на своё
\maketitle

\section*{Задача 1}

Покажите, что квадратная матрица $A$ коммутирует с $e^A$.

\section*{Задача 2}

\noindent Рассмотрим модель пружинного маятника на горизонтальной плоскости.

Его вектор состояния имеет вид

$x = \begin{pmatrix}
p \\
\dot{p} 
\end{pmatrix}$

Уравнение динамики системы в непрерывном случае таково:

$\dot{x} =  \begin{pmatrix}
0 & 1 \\
- \frac{k}{m} & 0 
\end{pmatrix} x$

Пусть симуляция этой системы производится как

$x_{k+1} = x_k + \Delta t \dot{x}_k$

Найдите, как изменяется полная энергия системы за один шаг симуляции.


\end{document}
